% RESULTADO_RECUPERADO: funtor dimensional integral, eq:cZ-internos,
% eq:longitud-ruta y definición de ell_0; composición con artículo III.
% La prueba reunida y su control algebraico son específicos de esta contribución.
\subsection{Vitesse constitutive et réalisation cinématique de l’horloge}
\label{v:sec:enlace-velocidad-constitutiva}

La propagation radiative utilise la section de vitesse produite par le
transport angulaire du même état HMT. L’identification est établie
avant le choix des unités de représentation. Soient $x=A_{\mathrm{rad}}$ et
$y=C^*_{\mathrm{rad}}$ les coordonnées angulaires engendrées dans la
section~\ref{v:ant:sec:ii-angulos}, avec les canaux
\eqref{v:ant:eq:ii-canales}, et soit
$\mathsf R$ l’involution autoadjointe des deux feuillets, avec projecteurs
de rang un $P_+$ et $P_-$. Dans le domaine $x>|y|$, le transport et ses
canaux sont
\begin{equation}
 T=\exp(-xI-y\mathsf R)=q_+P_++q_-P_-,
 \qquad q_+=e^{-x-y},\quad q_-=e^{-x+y},\quad 0<q_\pm<1.
 \label{v:eq:velocidad-transporte}
\end{equation}
La réponse des blocs temporels de longueurs $90$ et $120$ est
\begin{equation}
 \mathsf V=(I-T^{90})(I-T^{120})^{-1}
       =r_+P_++r_-P_-,\qquad
 r_\pm=\frac{1-q_\pm^{90}}{1-q_\pm^{120}}.
 \label{v:eq:velocidad-respuesta}
\end{equation}
La construction part des canaux hérités ; les symboles $x,y$ et
$q_\pm$ désignent leurs coordonnées et non des paramètres ajustés à une vitesse
cible.

\begin{proposition}[Identité des réalisations constitutive et cinématique]
\label{v:prop:identidad-velocidad}
Dans un même système de coordonnées dimensionnelles, la section $c_{\mathrm{int}}$ de l’horloge
et la section $c_{\mathrm{vac}}$ de la réponse constitutive coïncident :
\begin{equation}
 c_{\mathrm{int}}=c_{\mathrm{vac}}
   =\widehat c\,u_C,
 \qquad
 \widehat c=(r_+r_-)^{-1}
            =\exp(-\mathcal E),
 \label{v:eq:identidad-velocidad}
\end{equation}
où
\begin{equation}
 \mathcal E=
 \log\!\bigl[(1-q_+^{90})(1-q_-^{90})\bigr]
 -\log\!\bigl[(1-q_+^{120})(1-q_-^{120})\bigr].
 \label{v:eq:velocidad-lector-simetrico}
\end{equation}
L’égalité conserve l’échange des feuillets.
\end{proposition}
\begin{proof}
Les facteurs qui apparaissent dans les logarithmes sont positifs. Par conséquent,
\[
 \mathcal E=\log(r_+r_-)
            =\log\det\mathsf V,
 \qquad \exp(-\mathcal E)=(\det\mathsf V)^{-1}.
\]
Le déterminant est calculé ici dans la réalisation à deux feuillets de rang
un. La définition cinématique $c_{\mathrm{int}}=\exp(-\mathcal E)u_C$
et la définition constitutive $c_{\mathrm{vac}}=(r_+r_-)^{-1}u_C$
sont, par conséquent, la même section. L’échange $P_+\leftrightarrow
P_-$ permute $r_+$ et $r_-$ et conserve leur produit. Enfin,
$\det T=e^{-2x}$ correspond au transport élémentaire, tandis que
$\det\mathsf V=r_+r_-$ correspond à sa réponse temporelle ; les deux
déterminants ont des fonctions distinctes dans cette composition.
\end{proof}

\subsubsection{Longueur de trajectoire et durée élémentaire}

Le repère dimensionnel positif $(u_S,u_C,u_T,u_Q)$ induit la base de longueur $u_L=u_Cu_T$. Pour une trajectoire enrichie $\gamma$ de $n=|\gamma|$ transitions et à canal constant, les lecteurs additifs sont
\begin{equation}
 T_{\mathrm{int}}(\gamma)=n u_T,
 \qquad L_{\mathrm{int}}(\gamma)=n\widehat c\,u_L.
 \label{v:eq:velocidad-longitud-reloj}
\end{equation}
Son additivité se rapporte au dénombrement et à la réalisation de chaque transition :
les chemins enrichis conservent en outre l’orientation et la mémoire.

\begin{corollary}[Normalisation de la longueur élémentaire]
\label{v:cor:longitud-elemental-constitutiva}
Pour $n>0$, $L_{\mathrm{int}}(\gamma)/T_{\mathrm{int}}(\gamma)
=c_{\mathrm{vac}}$. Dans la normalisation $t_0=u_T$,
\begin{equation}
 \ell_0=c_{\mathrm{int}}t_0=\widehat c\,u_L.
 \label{v:eq:longitud-elemental-constitutiva}
\end{equation}
Dans un système de coordonnées temporelles général $t_0=\tau_0u_T$, avec $\tau_0>0$,
l’expression correspondante est $\ell_0=\tau_0\widehat c\,u_L$.
\end{corollary}
\begin{proof}
La division de \eqref{v:eq:velocidad-longitud-reloj} simplifie l’entier
$n$ et utilise $u_L/u_T=u_C$. La formule élémentaire s’obtient
en multipliant la section de vitesse par $t_0$.
\end{proof}

La généralisation à un canal dépendant de l’arête conserve la formule
additive $L_{\mathrm{int}}(\gamma)=\sum_{e\in\gamma}\widehat c(e)u_L$.
Son quotient par $nu_T$ est la moyenne des vitesses de ces arêtes.
La réalisation radiative homogène emploie le canal constant déclaré
dans \eqref{v:eq:velocidad-longitud-reloj}.

\subsubsection{Changement de coordinatisation et normalisation adaptée}

\begin{proposition}[Covariance de l’identification]
\label{v:prop:covariancia-identidad-velocidad}
Soient $u_C'=d u_C$ et $u_T'=\eta u_T$, avec $d,\eta>0$.
La même vitesse et la même durée élémentaire ont pour coordonnées
\begin{equation}
 \widehat c'=d^{-1}\widehat c,
 \qquad \tau_0'=\eta^{-1}\tau_0,
 \qquad u_L'=d\eta u_L.
 \label{v:eq:velocidad-cambio-carta}
\end{equation}
En particulier, si $c_V=\widehat c_Vu_C^V$ et
$u_C^V=d u_C^{III}$, l’égalité des sections équivaut exactement à
$d\widehat c_V=\widehat c_{III}$.
\end{proposition}
\begin{proof}
On exprime chaque section dans les deux bases. On obtient
$\widehat c'u_C'=\widehat c u_C$ et
$\tau_0'u_T'=\tau_0u_T$. Leur produit satisfait
\[
 \tau_0'\widehat c' u_L'
 =\frac{\tau_0}{\eta}\frac{\widehat c}{d}
       d\eta u_L
 =\tau_0\widehat c u_L=\ell_0.
\]
La dernière équivalence résulte de la substitution de la relation entre les bases.
\end{proof}

Le choix adapté $d=\widehat c$ donne $u_C'=c_{\mathrm{vac}}$ et
$\widehat c'=1$. Ce choix exprime la section déjà construite dans une
base qui lui est adaptée. Le coefficient spectral dans le système de coordonnées interne reste
$(r_+r_-)^{-1}$ ; l’inversion spectrale utilise ce système de coordonnées interne.
Par exemple, en maintenant les bases d’action et de charge, les coordonnées
constitutives se transforment comme
\[
 \widehat\mu'=d\widehat\mu,
 \qquad \widehat\varepsilon'=d\widehat\varepsilon,
 \qquad
 \widehat\mu'\widehat\varepsilon'=(\widehat c')^{-2}.
\]
Dans le système de coordonnées adapté, $\widehat\mu'=r_-/r_+$ et
$\widehat\varepsilon'=r_+/r_-$, car les coordonnées internes sont
$\widehat\mu=r_-^2$ et $\widehat\varepsilon=r_+^2$.

\subsubsection{Moment orienté de Catalan et pont fonctionnel}

La réponse constitutive ne reçoit pas la constante de Catalan comme un
nombre externe. Le quart de tour dodécaphasique produit la collection de
moments orientés
\begin{equation}
 \mathcal C_2(s)=\sum_{k\geq0}
 \frac{(-1)^ks^{2k+1}}{(2k+1)^2},\qquad
 \mathcal D=s\frac{\dd}{\dd s},\qquad 0<s<1.
 \label{v:eq:momento-catalan}
\end{equation}
La convergence absolue sur les compacts permet de dériver terme à terme et
donne
\begin{equation}
 \mathcal D^2\mathcal C_2(s)=\frac{s}{1+s^2},\qquad
 f(s)=\frac{1+s+s^2}{s(1+s)}\mathcal D^2\mathcal C_2(s)
 =\frac{1-s^3}{1-s^4}.
 \label{v:eq:puente-catalan-constitutivo}
\end{equation}
Avec $s_\pm=q_\pm^{30}$, les incidences $90=3\cdot30$ et
$120=4\cdot30$ produisent exactement
\begin{equation}
 r_\pm=f(s_\pm),\qquad
 \mathcal C_2(1)=G_{\mathrm{Cat}}.
 \label{v:eq:catalan-canales}
\end{equation}
La valeur extrême est la représentation de la collection ; la collection fonctionnelle,
et non cette seule valeur, gouverne les deux canaux.

Le pont est réversible dans la classe régulière fixée. Étant donnée $f$, la
condition $\mathcal C_2(s)\to0$ lorsque $s\downarrow0$ sélectionne
de manière unique
\begin{equation}
 \mathcal C_2(s)=\int_0^s\log\frac{s}{t}\,
 \frac{1+t}{1+t+t^2}f(t)\,\dd t.
 \label{v:eq:restitucion-funcional-catalan}
\end{equation}
En effet, cette intégrale satisfait
$\mathcal D^2\mathcal C_2=s(1+s)f/(1+s+s^2)$ ; pour
\eqref{v:eq:puente-catalan-constitutivo}, l’intégrande se réduit à
$\log(s/t)/(1+t^2)$. Deux solutions homogènes diffèrent de
$a\log s+b$, et la régularité nulle à l’origine impose $a=b=0$.

La même collection dilogarithmique organise un potentiel conjoint. Soit
$\mathcal D_2(z)=\sum_{n\geq1}z^n/n^2$ et
\begin{equation}
 \mathscr F(x,y)=\sum_{\sigma=\pm}
 \left[
 \frac{\mathcal D_2(q_\sigma^{120})}{120}
 -\frac{\mathcal D_2(q_\sigma^{90})}{90}
 \right].
 \label{v:eq:potencial-constitutivo-catalan}
\end{equation}
Comme $\dd\mathcal D_2(z)/\dd\log z=-\log(1-z)$ et
$q_\pm=e^{-x\mp y}$,
\begin{equation}
 \partial_x\mathscr F=\log\widehat c,
 \qquad
 \partial_y\mathscr F=\log\widehat Z,
 \qquad \widehat Z=\frac{r_-}{r_+}.
 \label{v:eq:gradiente-constitutivo-catalan}
\end{equation}
Ainsi, la vitesse et l’impédance sont les deux dérivées conjuguées d’un même
potentiel, avec une constante additive fixée par $\mathscr F\to0$ lorsque
$q_\pm\to0$.

\subsubsection{Restitution conjointe au moyen de la vitesse et de Barbero--Immirzi}

Soit $\psi=f^{-1}$ et définissons
\begin{equation}
 \mathcal H_B(s)=\frac{12s^3}{1-s^9}
 -\frac{s^4}{1-s^{12}}-\frac{(\log s)^2}{20\pi}.
 \label{v:eq:funcional-barbero-espectral}
\end{equation}
La formule de \eqref{v:ant:eq:gamma}, avec
\eqref{v:eq:catalan-canales}, se factorise exactement comme
\begin{equation}
 \gamma_{\HMT}=\mathcal H_B(\psi(r_-))
 -\mathcal H_B(\psi(r_+)).
 \label{v:eq:factorizacion-barbero-espectral}
\end{equation}

\begin{theorem}[Restitution conjointe des canaux étiquetés]
\label{v:thm:restitucion-c-gamma}
Dans la collection constitutive normalisée, l’application
\begin{equation}
 \Psi:(3/4,1)^2\longrightarrow(1,16/9)\times\mathbb R,
 \qquad
 (r_+,r_-)\longmapsto(\widehat c,\gamma_{\HMT})
 \label{v:eq:mapa-c-gamma}
\end{equation}
est un difféomorphisme. Par conséquent, $(\widehat c,\gamma_{\HMT})$ restitue
les deux valeurs propres étiquetées ; les projecteurs $P_\pm$ restent des
données de la représentation initiale.
\end{theorem}
\begin{proof}
L’expression régulière $f(s)=(1+s+s^2)/(1+s+s^2+s^3)$ satisfait
\[
 f'(s)=-\frac{s^2(s^2+2s+3)}{(1+s+s^2+s^3)^2}<0,
\]
avec $f(0^+)=1$ et $f(1^-)=3/4$. En outre,
\begin{equation}
 \mathcal H_B'(s)=
 \frac{36s^2(1+2s^9)}{(1-s^9)^2}
 -\frac{4s^3(1+2s^{12})}{(1-s^{12})^2}
 -\frac{\log s}{10\pi s}>0.
 \label{v:eq:monotonia-hb}
\end{equation}
Le quotient du premier terme positif par le second avant sa soustraction
est
\[
 \frac9s\frac{1+2s^9}{1+2s^{12}}
 \left(\frac{1-s^{12}}{1-s^9}\right)^2>9,
\]
et le dernier terme de la somme est aussi positif. Ainsi,
$F_B=\mathcal H_B\circ\psi$ est strictement décroissante, avec limites
$F_B(3/4^+)=+\infty$ et $F_B(1^-)=-\infty$.

$c_0=\widehat c$ étant fixé et $z=\log\widehat Z$ étant écrit, on a
$r_\pm=c_0^{-1/2}e^{\mp z/2}$ et
\[
 |z|<a(c_0):=\min\!\left\{\log c_0,
 \log\frac{16}{9c_0}\right\}.
\]
La fonction $\Gamma_{c_0}(z)=F_B(r_-)-F_B(r_+)$ satisfait
\[
 \Gamma_{c_0}'(z)=\frac12
 \bigl[r_-F_B'(r_-)+r_+F_B'(r_+)\bigr]<0
\]
et parcourt $\mathbb R$ entre les deux extrémités admissibles. Il existe une
solution unique pour chaque $\gamma_{\HMT}$ ; la dérivée non nulle et le théorème
des fonctions implicites fournissent l’inverse lisse global.
\end{proof}

Sur l’image produite par le générateur HMT, la restitution conserve
l’ordre causal. Une fois $r_\pm$ récupérés, on obtient
$s_\pm=\psi(r_\pm)$ et
\begin{equation}
 A_{\deg}=-\frac3\pi\log(s_+s_-),\qquad
 C^*_{\deg}=\frac3\pi\log\frac{s_-}{s_+},
 \label{v:eq:restitucion-angular-constitutiva}
\end{equation}
d’où, dans l’orientation positive déjà engendrée,
$\alpha_{\HMT}=A_{\deg}/1000$ et
$\eta_{\mathrm{ret}}=C^*_{\deg}/2-\alpha_{\HMT}$. Ce parcours
reconstruit des représentations de l’état ; il ne transforme ni
$(\widehat c,\gamma_{\HMT})$ ni des valeurs conventionnelles en entrées du
générateur.

\subsubsection{Transport vers le secteur radiatif}

La dispersion de la réalisation radiative homogène s’écrit
$\omega_{\symbf{k}}=c_{\mathrm{vac}}|\symbf{k}|$.
Avec $h=2\pi\hbar$ et $\beta=(k_B\Theta)^{-1}$, sa densité spectrale
et son coefficient de Stefan--Boltzmann admettent les expressions
\begin{equation}
 u_\omega=
 \frac{\hbar\omega^3}{\pi^2(\widehat c u_C)^3}
 \frac{1}{e^{\beta\hbar\omega}-1},
 \qquad
 \sigma=\frac{\pi^2k_B^4}{60\hbar^3(\widehat c u_C)^2}.
 \label{v:eq:radiacion-velocidad-constitutiva}
\end{equation}
Ces substitutions transportent exactement la même section de
vitesse. Les conditions de la réalisation radiative sont conservées :
secteur bosonique libre tridimensionnel, deux polarisations transversales,
potentiel chimique nul et température positive. L’identité
\eqref{v:eq:identidad-velocidad} établit son antécédent constitutif ;
le dénombrement des modes et la limite thermodynamique appartiennent aux preuves
spécifiques de cette réalisation.

\paragraph{Normalisation sous amplification.}
Si le raffinement agit par $\mathsf V_m=\mathsf V\otimes I_{9^m}$,
la trace normalisée $\tau_m=\operatorname{Tr}/(2\cdot9^m)$ satisfait
\[
 \exp[-2\tau_m(\log\mathsf V_m)]
 =\exp[-\log r_+-\log r_-]=\widehat c.
\]
En effet, chaque valeur propre apparaît $9^m$ fois. La multiplicité est
simplifiée dans la trace normalisée. Le déterminant ordinaire agrandi est
$(r_+r_-)^{9^m}$ : la normalisation à deux feuillets, ou de manière équivalente la
trace normalisée, préserve le lecteur constitutif pendant ce
raffinement. Cette amplification de la représentation conserve le
transport enrichi qui en est à l’origine.
